\clearpage\thispagestyle{empty}
\pdfbookmark[0]{Sources for Determining Longitude by Chronometer}{chronometric1826-sources}
\begin{center}\textbf{Sources for Determining Longitude by Chronometer}\end{center}
\noindent This AI translation follows \textit{Chronometrische Längenbestimmungen} in \href{https://gdz.sub.uni-goettingen.de/id/PPN236007467}{\textit{Werke}, vol.~VI} (Göttingen, 1874), pp.~455--459. The first printing in \textit{Astronomische Nachrichten}, vol.~V, no.~110, calls the paper an extract of a letter from Gauss to the editor. The Bavarian State Library copy gives the text in cols.~\href{https://api.digitale-sammlungen.de/iiif/image/v2/bsb10049430_00124/full/1500,/0/default.jpg}{227--228}, \href{https://api.digitale-sammlungen.de/iiif/image/v2/bsb10049430_00125/full/1500,/0/default.jpg}{229--230}, \href{https://api.digitale-sammlungen.de/iiif/image/v2/bsb10049430_00126/full/1500,/0/default.jpg}{231--232} and \href{https://api.digitale-sammlungen.de/iiif/image/v2/bsb10049430_00127/full/1500,/0/default.jpg}{233--234}. The collected edition dates the item November 1826. The bound journal volume is dated 1827.
\par\smallskip\noindent The collected edition includes Ernst Schering’s note identifying Dr. Tiarks’s observations for the English Admiralty expedition of 1824. Its text repeats the number 4, has a raised stroke after the $S$ under the first radical, and gives the first Greenwich comparison in June. The journal printing repeats 3, prints an unprimed $S$ under the first radical and gives January in the prose while its table has June. The journal also names the offset of the imaginary chronometer in the Bremen calculation and closes with Gauss’s signature. The translation follows the collected text and table arrangement. Translation and typesetting produced by AI.
\clearpage
